\section{Conclusions}\label{sec:conclusions}

On the open biomedical citation graph, papers that were later retracted, including those retracted for misconduct and
paper-mill production, have a small mean CD$_5$ difference from comparable papers, and no unadjusted excess in the most
disruptive tier in the pre-specified five-year analysis. The post hoc three-year window leans toward a modest excess,
though this weakens under author-name exclusions. Positive scores are nevertheless more common, a sign shift that should not be read as complete
agreement between the groups. These conclusions concern papers retracted only after their scoring window.
The exploratory DI$_{nok}$ comparison shows what the index's largest term, $n_k$, compresses:
papers that cite a retracted paper are less likely to also cite its references. This difference persists after
removing notices, comments/letters and retracted citing papers, adjusting reference and citation counts, and
excluding the ten highest-retraction journals. A variant without $n_k$ therefore scores these late-retracted papers
somewhat higher than their controls. Scores that include retraction years behave differently again, through
notice records and a widening gap; scoring windows should end before the first notice, with notice records removed.
The analysis can be repeated with later snapshots of the open graph, which will add
the large cohorts of paper-mill papers published after 2019.
